\snote{robust}

This note describes in detail how the robustness of the selected subsets to the size of the parameter changes was assessed (Fig.~4 of the main text) and gives, for every admissible problem, the full distribution of the fitting error over the random scenarios, not only its median.

\subsection*{Random scenarios}
For a problem with $p$ estimated parameters, fifteen random directions $u^{(1)},\dots,u^{(15)}\in[-1,1]^p$ are drawn once, with independent components $u^{(k)}_j\sim\mathcal U(-1,1)$, from a pseudo-random generator with fixed seed (NumPy \texttt{default\_rng(42)}). For a perturbation size $\varepsilon$, scenario $k$ changes all parameters simultaneously:
\[ \theta^{(k)}(\varepsilon)=\theta_0\odot\bigl(\mathbf 1+\varepsilon\,u^{(k)}\bigr),\qquad \theta^{(k)}_j(\varepsilon)=\theta_{0,j}\,\bigl(1+\varepsilon\,u^{(k)}_j\bigr),\quad j=1,\dots,p, \]
so that every parameter is multiplied by a factor between $1-\varepsilon$ and $1+\varepsilon$. The same fifteen directions are used at every size $\varepsilon\in\{0.01,0.05,0.10,0.20,0.30,0.40,0.50\}$; only $\varepsilon$ changes (paired design). Differences between sizes are therefore due to the size of the change and not to different random draws, and the scenarios at $\varepsilon=0.05$ are exactly those used to select the subset. The directions of every problem are stored in the result files (\texttt{certificados/resultados/patrones\_escenarios/<problem>.json}, key \texttt{u}), and the script checks that they are identical at every size.

\subsection*{Errors in each scenario}
Let $y(\theta)\in\mathbb R^m$ be the vector of noise-scaled measured outputs (main text, Methods) and $S$ the selected biological parameters together with the calibration parameters $\Theta_C$. For scenario $k$,
\[ \Delta y^{(k)}_{\mathrm{full}}=y\bigl(\theta^{(k)}\bigr)-y(\theta_0),\qquad \Delta y^{(k)}_{S}=y\bigl(\theta^{(k)}_S\bigr)-y(\theta_0),\qquad \bigl(\theta^{(k)}_S\bigr)_j=\begin{cases}\theta^{(k)}_j,& j\in S\\ \theta_{0,j},& j\notin S,\end{cases} \]
that is, $\Delta y_S$ is the output change when only the retained parameters take their scenario values. The error without re-estimation and the cosine are
\[ e^{(k)}_{\mathrm{rel}}=\frac{\|\Delta y^{(k)}_{\mathrm{full}}-\Delta y^{(k)}_{S}\|}{\|\Delta y^{(k)}_{\mathrm{full}}\|},\qquad \cos\Delta^{(k)}=\frac{\langle\Delta y^{(k)}_{\mathrm{full}},\Delta y^{(k)}_{S}\rangle}{\|\Delta y^{(k)}_{\mathrm{full}}\|\,\|\Delta y^{(k)}_{S}\|}. \]
The error after re-estimation is the smallest relative distance between the reduced model and the output of the full model in the scenario, when the retained parameters are free and all others are fixed at $\theta_0$:
\[ e^{(k)}_{\mathrm{fit}}=\min\Bigl\{e^{(k)}_{\mathrm{rel}},\ \min_{v\in B}\frac{\|y(\theta_0[S\leftarrow v])-y(\theta^{(k)})\|}{\|\Delta y^{(k)}_{\mathrm{full}}\|}\Bigr\},\qquad B=\prod_{j\in S}\Bigl[\tfrac{|\theta_{0,j}|}{10},\,10\,|\theta_{0,j}|\Bigr], \]
where $\theta_0[S\leftarrow v]$ is $\theta_0$ with the components in $S$ replaced by $v$. The inner minimisation is a bounded nonlinear least-squares problem in $\log|v|$ (trust-region reflective method, at most $100(|S|+1)$ function evaluations, finite-difference step $10^{-3}$), started from the scenario values (clipped to $B$) and from $\theta_0$; the smaller result is kept. Since $v=\theta^{(k)}_S$ is admissible when it lies in $B$, $e_{\mathrm{fit}}\le e_{\mathrm{rel}}$ always holds (this is also proved in Lean, step~5). For the problems whose simulations are too slow for repeated nonlinear fits (Alkan, Bachmann, Chen, Isensee, Lang, Lucarelli, Raimundez), the first-order version is used,
\[ e^{(k)}_{\mathrm{fit,lin}}=\min\Bigl\{e^{(k)}_{\mathrm{rel}},\ \min_{c\in\mathbb R^{|S|}}\frac{\|\Delta y^{(k)}_{\mathrm{full}}-J_S\,c\|}{\|\Delta y^{(k)}_{\mathrm{full}}\|}\Bigr\}, \]
with $J_S$ the columns of the relative sensitivity matrix for $S$ (central differences, step 1\%); it differs from the exact value by a second-order term (Lean, step~12).
A scenario is discarded when the simulation fails or when the full-model change is negligible, $\|\Delta y^{(k)}_{\mathrm{full}}\|<\max(10^{-10}\|y(\theta_0)\|,10^{-6})$, because the relative errors are then undefined; the number $n\le15$ of valid scenarios is reported for every problem and size.

\subsection*{Summaries over the scenarios}
For each problem and size $\varepsilon$, with $e^{(1)},\dots,e^{(n)}$ the values of $e_{\mathrm{fit}}$ in the valid scenarios:
\begin{itemize}\itemsep0pt
\item median $\tilde e_{\mathrm{fit}}$, first and third quartiles $Q_1,Q_3$ (interquartile range $Q_3-Q_1$), and the 5th and 95th percentiles $P_5,P_{95}$, computed by linear interpolation between order statistics (with $n=15$, $P_5$ and $P_{95}$ lie between the two smallest and the two largest values);
\item maximum $\max_k e^{(k)}_{\mathrm{fit}}$, the worst scenario;
\item pass rate, the fraction of scenarios in which the reduced model reproduces the full one,
\[ \pi(\varepsilon)=\frac{\#\{k:\ e^{(k)}_{\mathrm{fit}}\le\tau_e\}}{n},\qquad \tau_e=\sqrt{1-0.9^2}=0.436. \]
\end{itemize}
Admissibility at $\pm5\%$ and the robustness patterns of the main text use the median: pattern A, $\tilde e_{\mathrm{fit}}\le\tau_e$ at every size; B, admissible up to at least $\pm20\%$ and not beyond; C, irregular (fails and is admissible again at a larger size); D, not admissible at $\pm10\%$ or less. The median requires that at least half of the scenarios are reproduced; the pass rate and the maximum show how many, and how badly the others fail.

\subsection*{Results}
\textit{Per-scenario results are being computed for 31 problems; this section is generated automatically when they are complete.}
